% !TeX root = ../papers/nc-p3.tex
\section{七次元の余因子空間と二つの有理正規四次曲線の共通三割平面}
\label{sec:seven-cofactor-trisecant-grassmannian}
\begin{theorem}\label{thm:seven-cofactor-trisecant-grassmannian}
有限Frobenius窓の異なる反対根について、$D=V(\det N)\subset\mathbb P^3$ が
整な八次曲面であるとする。六次余因子の線型包の次元 $m$ は7以上である。
$m=8$ なら一般核点の像の閉包は二つの二次式の完全交叉である整四次曲線、
$m=7$ なら有理正規三次曲線である。後者の行列は入出力基底を除き次の形に限られる。
$S_i$ を二変数の $i$ 次形式、$T\in\mathrm{GL}(S_4)$ として
\[
       N(x)=\operatorname{Bez}\big(T(sx),T(tx)\big),\qquad x\in S_3 .
\]
$\operatorname{Bez}$ は二つの四次式のBezout行列である。
配置 $X$ の最初と反対の点は、二つの有理正規四次曲線の共通三割平面を与える二つの
長さ3因子の座標になる。整な八次の場合にこの二曲線は交わらない。
行列としての上の表示の同値は、$[T]\in\mathrm{PGL}(S_4)$ の左右からの
$\operatorname{Sym}^4\mathrm{PGL}_2$ による二重剰余で決まる。
\end{theorem}

\begin{proof}[余因子と核曲線]
八次式は既約なので $D$ の一般点で $N$ の階数は3である。
核直線の像の閉包を $Z$ とすれば、余因子への定数線型関係はちょうど $Z$ の二次式である。
実際余因子は $D$ 上で核点の平方のスカラー倍であり、$D$ 上消える六次式は恒等零である。
$Z$ は非退化である。ある一次形式 $l$ で核を全て消せれば
$(\operatorname{adj}N)l$ の全成分が $D$ 上で消えて恒等零となり、一般可逆性に反する。
二つ以上の二次式があれば、それらは整で非退化な $Z$ を含むので可約ではなく、
独立な二式は固有に交わる。$Z$ は三次または四次の非退化整曲線となる。
最小次数3の場合は正規化で次数3の線束が独立な四切断を持つので有理正規三次曲線であり、
四次なら次数の等号と完全交叉の純性より二式の全交叉である。
前者の二次式は三次元、後者のものは二次元である。これで $m$ の主張が従う。

$m=7$ の核曲線を $(1,t,t^2,t^3)$ の閉包に取ると、$\operatorname{adj}N(x)$ は
反対角線に沿う成分が等しいHankel行列である。
可逆Hankel行列を逆にした像の閉包は、以下の固定した線型同型のもとで
$\operatorname{Gr}(2,S_4)\subset\mathbb P(\wedge^2 S_4)$ である。
交代な二変数四次式テンソルを二点の行列式で割ると、対称な三次式テンソルになる。
これは線型同型で、分解可能な $f\wedge g$ に対する像が $\operatorname{Bez}(f,g)$ である。
一般に $f$ が異なる四根を持ち $g$ と共通根がないとき、各根での三次評価の四基底で
Bezout対合は対角かつ可逆である。逆行列は各根の評価ベクトルの平方の線型結合で
Hankelとなる。両閉包は整で次元6（Grassmann多様体とHankel線型空間の双有理像）なので一致する。
\end{proof}

\begin{proof}[Grassmann多様体内の二次Veronese]
従って $x\mapsto[N(x)]$ は三空間の二次Veroneseとして
$\operatorname{Gr}(2,5)$ に入る。普遍部分束の双対を引き戻して $F$ とすると
五切断が生成する階数2の束で、$\det F=\mathcal O(2)$、
五切断の外積は全二次式を張る。一般切断の零点が空なら $F=\mathcal O\oplus\mathcal O(2)$
であるが、この場合五切断のうち四本を第二因子に入れられ、外積は高々四次元である。
空でなければ一般零点は滑らかな曲線 $C$ で、随伴は $\omega_C=\mathcal O_C(-2)$。
各成分の次数と種数から $C$ は互いに素な直線の和である。
$\mathcal I_C(2)$ は大域生成なので三本以上はない。三本のねじれの位置の直線を含む
二次式は一次元だけである。二本を座標直線に取ると二次式は二組の変数の双一次式で、
第三の直線は両二次元ベクトル空間の間の同型のグラフなので交代条件が残るためである。

二直線なら $h^0(F)\le1+h^0(\mathcal I_C(2))=5$ であり、
$h^0(\operatorname{Sym}^2 F)\le5+h^0(\mathcal I_C^2(4))=14$ となる。
後者では零点を定める切断での $F\to\operatorname{Sym}^2F$ の余核が
$\mathcal I_C^2(4)$ となることを使った。二生成正則列の局所表示でその二次の積の
二つの関係を展開すればよい。ねじれた二直線での二重消滅四次式は二組の変数について
二重次数 $(2,2)$ の九次元である。従って五空間の非零二次式 $Q$ が全てのパラメータ直線で消える。
一方、階数1の対称行列 $yy^{\mathsf t}$ に対応するGrassmann点は二平面 $yS_1$ である。
二点行列式と $y$ の両評価の積が $(sy)\wedge(ty)$ だからである。
$X$ の反対射影上で $Q(sy),Q(ty)$ と偏極 $Q(sy,ty)$ は全て零となる。
これらは $y$ の二次式でありその評価は単射なので、全 $y$ で消える。
全ての二項四次式は一次因子を持つから二平面 $yS_1$ の和集合は四次式空間全体であり矛盾する。

一線の場合、$0\to\mathcal O\to F\to\mathcal I_C(2)\to0$ は
$F\simeq\mathcal O(1)^2$ を与える。直線の二生成分解からこの拡張空間は一次元で、
非分裂拡張を $\mathcal O(1)^2$ が与え、分裂拡張は局所自由でない。
よって入力 $x$ の直線族は二つの線型写像 $A,B:k^4\to k^5$ による
$\langle Ax,Bx\rangle$ である。各 $x\ne0$ で二列は独立なので
$uA+vB$ は全ての $[u:v]$ で単射である。
$\mathbb P^1$ 上の $\mathcal O(-1)^4\hookrightarrow\mathcal O^5$ の余核は
$\mathcal O(4)$ であり、$H^0(\mathcal O(-1))=H^1(\mathcal O(-1))=0$ より
五切断は四次式の基底となる。この全評価の核は、点で消える一次形式と任意の三次式との積
$\mathcal O(-1)\otimes S_3$ である。定数基底変更で $(A,B)$ は
$(x\mapsto T(sx),x\mapsto T(tx))$ となる。Plücker座標の共通定数は一方の生成列に吸収できる。
\end{proof}

\begin{proof}[三割平面と行列表示の同値]
$\mathbb P(S_4^*)$ の四次Veronese曲線を $R$ とする。
$yS_1$ の零化平面は $y$ の三根の $R$ 上の張る平面で、重根なら対応する長さ3因子を使う。
実際その三点条件を満たす四次式はちょうど $y$ の一次式倍である。
従って二つの三空間は $R$ と $T^{-\mathsf t}R$ の三割平面族であり、$X$ はその共通部分に入る。
両曲線が共有点を持てば、その評価で $T(sx),T(tx)$ がともに消える条件は
$x$ の一つの一次式の消滅に等しくなる。元の評価が $T$ の引戻しでも一点評価だからである。
従ってその平面はBezout行列式に入る（共通根があればその評価ベクトルが核に入る）。
整な八次という仮定に反する。

核三次曲線からの出力基底の自由度はその三次対称冪の $\mathrm{PGL}_2$ であり、
Bezout同型ではこれは四次対称冪の作用になる。入力の三空間の基底は族のパラメータだけを変える。
$\operatorname{Gr}(2,S_4)$ 上で標準族 $\{yS_1\}$ を保つ $\mathrm{PGL}(S_4)$ の部分群は
四次対称冪の $\mathrm{PGL}_2$ に等しい。実際 $[f]\in\mathbb P(S_4)$ を通る族の直線数は
$f$ の異なる一次因子の数であり、ちょうど一本の点は純四乗曲線である。
その安定群は曲線上の射影直線の自己同型で尽きる（曲線は射影枠を含む）。
これで二つの表示が同じ基底軌道に入るときの必要性が従う。
逆に両側のその基底変更は実際に行列を合同と入力射影変換で移す。
これは標識付きの行列表示の判定であり、追加の隣接座標を選ぶには関係の検査が要る。
\end{proof}

\begin{corollary}\label{cor:trisecant-middle-edge}
上の $T$ に対して共通三割平面は、中間座標 $[l:m]\in\mathbb P(S_1\oplus S_1)$ を使い、
六つの双一次式
\[
                  l\,T(sx)+m\,T(tx)=0\quad\text{in } S_5
\]
で入射表示できる。中間座標への退化スキーム（六行四列行列の最大小行列式）は常に有限で長さ20である。
その全点の評価階数が3なら、共通三割平面スキーム自体も長さ20で、既に含まれている
$X$ を復元する。この場合はこの中間座標を両奇位置に取る候補の標準正則性を有限判定できる。
\end{corollary}
\begin{proof}
$T(sx),T(tx)$ は $x\ne0$ で独立である。二本の線型係数の関係を持てば
$l,m$ は互いに比例せず、両四次式が同じ三次式 $y$ の一次式倍になる。
逆も成立し、係数関係の直線は一意である。これはスキームとしても、$l,m$ の係数行列の
単元行列式を使い、二本の一次式によるKoszul関係を解けば成立する。
共通三割平面からは二平面の二つの基底の間の可逆行列により、この中間座標を復元できる。
従って中間射影の像は射影三空間の行列式二次曲面を避ける。射影的な像がこのアフィン開集合に入るので
零次元である。最大小行列式の高さは高々3なので空でもなく、期待余次元の
六行四列の一次行列の分解から長さは $\binom63=20$ である。
階数3なら各点で核直線をスキームとして解ける。逆方向の $T^{-1}$ による中間座標は
同じ基底変更行列の逆であり、二行二列の射影逆は余因子による線型変換である。
従って両奇座標が射影同値となる候補を二辺から再構成し、有限Frobenius検査と
既出の一交換対称の有限正則判定を適用できる。
\end{proof}

\begin{proposition}\label{prop:trisecant-common-quadric}
上の中間座標による隣接二辺で、六式が両側で独立なら、三次交叉の次元は
$3+h^0(\mathcal I_{R\cup T^{-\mathsf t}R}(2))$ である。
特にこの二辺が標準延長を持つには二つの四次曲線が共通の二次超曲面をただ一つ持つことが必要である。
一般の $T$ に対しては三次交叉は三次元で、この二辺による延長はない。
\end{proposition}
\begin{proof}
$s,t$ を $u_0,u_1$ と書く。中間文字は $l,m$ の四係数で、
二番目の辺では添字を転置した四係数を使う。
実際 $l=l_0 s+l_1t$、$m=m_0s+m_1t$ とし、$T(sx)=my,T(tx)=-ly$ とすると、
逆の二四次式の一次シジジーは $(l_0s+m_0t,l_1s+m_1t)$ である。
従って交叉は線型写像 $a,b:S_3\to S_5^*$ による次の等式の解である。
\[
     a(y)(u_j T(u_i x))=b(x)(u_i T^{-1}(u_j y))
              \qquad (i,j=0,1).
\]
左辺は $u_j y$ に関する線型関係を右辺により保ち、右辺は $u_i x$ に関する線型関係を
左辺により保つ。よって $k=u_i x,h=u_j y$ の間の双線型形式 $B(k,h)$ に降りる。
逆に $B$ が $a$ を与えるには、$u_0S_4\cap u_1S_4=u_0u_1S_3$ における二表示が一致すればよい。
これは $G(k,h)=B(T^{-1}k,h)$ が
$G(u_0 w,u_1y)=G(u_1w,u_0y)$ を満たすこと、すなわち五行Hankel形式であることに等しい。
同様に $b$ が存在する条件は $B(k,Th)$ がHankelであることで、両写像は存在すれば一意となる。
従って解はHankel九空間 $\mathcal H$ に対する $G\in\mathcal H$、$T^{\mathsf t}GT\in\mathcal H$
である。この九空間は曲線 $R$ の評価ベクトルの平方の包であり、その零化空間は $I_R(2)$。
二つの九空間の和の零化空間を十五次元の対称二次空間で取れば次元公式を得る。
一般値については付属の有理整数行列で二つの九空間の包が十五次元、両側の六式が独立、
元の二つの二十四次元空間の包が四十五次元であることを有理数上で確認できる。
いずれも最大階数の開条件である。
\end{proof}
